\section{Задание 3. Матрица линейного оператора}

1. Выяснить, какие из следующих операторов $\varphi$ пространства $\mathbb{R}^3$, заданных путём задания координат вектора $\varphi(x)$ как функций координат вектора
\[
x=(x_1,x_2,x_3)^T,
\]
являются линейными.

2. Если оператор линейный, найдите его матрицу в стандартном базисе $e_1,e_2,e_3$ и в базисе $e_1',e_2',e_3'$, где
\[
\begin{gathered}
e_1' = e_1 - e_2 + e_3,\\
e_2' = -e_1 + e_2 - 2e_3,\\
e_3' = -e_1 + 2e_2 + e_3.
\end{gathered}
\]

\[
17.\quad
\begin{aligned}
\text{а)}\ \varphi(x) &= (x_1 + x_3,\,-x_2 + 2x_3,\,3x_1 - x_2 + x_3)^T;\\
\text{б)}\ \varphi(x) &= (x_1 + x_3,\,-x_2 + 2x_3,\,3x_1 - x_2 + 1)^T;\\
\text{в)}\ \varphi(x) &= (x_1 + x_3,\,-x_2 + 2x_3,\,3x_1 - x_2^2 + x_3)^T.
\end{aligned}
\]

\textbf{Решение.}

\textbf{1. Выясним, какие из операторов являются линейными.}

Оператор $\varphi$ линеен, если для любых $x,y\in\mathbb{R}^3$ и $\lambda\in\mathbb{R}$ выполняются условия:
\[
\varphi(x+y)=\varphi(x)+\varphi(y),
\qquad
\varphi(\lambda x)=\lambda \varphi(x).
\]

\textbf{а)} Проверим, является ли линейным оператор
\[
\varphi(x)=
\begin{pmatrix}
x_1+x_3\\
-x_2+2x_3\\
3x_1-x_2+x_3
\end{pmatrix}.
\]

Пусть
\[
x=(x_1,x_2,x_3)^T,\qquad y=(y_1,y_2,y_3)^T.
\]

Тогда
\[
x+y=(x_1+y_1,\;x_2+y_2,\;x_3+y_3)^T,
\]
и
\[
\varphi(x+y)=
\begin{pmatrix}
(x_1+y_1)+(x_3+y_3)\\
-(x_2+y_2)+2(x_3+y_3)\\
3(x_1+y_1)-(x_2+y_2)+(x_3+y_3)
\end{pmatrix}
\]
\[
=
\begin{pmatrix}
x_1+x_3+y_1+y_3\\
-x_2+2x_3-y_2+2y_3\\
3x_1-x_2+x_3+3y_1-y_2+y_3
\end{pmatrix}.
\]

С другой стороны,
\[
\varphi(x)+\varphi(y)=
\begin{pmatrix}
x_1+x_3\\
-x_2+2x_3\\
3x_1-x_2+x_3
\end{pmatrix}
+
\begin{pmatrix}
y_1+y_3\\
-y_2+2y_3\\
3y_1-y_2+y_3
\end{pmatrix}
\]
\[
=
\begin{pmatrix}
x_1+x_3+y_1+y_3\\
-x_2+2x_3-y_2+2y_3\\
3x_1-x_2+x_3+3y_1-y_2+y_3
\end{pmatrix}.
\]

Следовательно,
\[
\varphi(x+y)=\varphi(x)+\varphi(y).
\]

Проверим однородность:
\[
\varphi(\lambda x)=
\begin{pmatrix}
\lambda x_1+\lambda x_3\\
-\lambda x_2+2\lambda x_3\\
3\lambda x_1-\lambda x_2+\lambda x_3
\end{pmatrix}
=
\lambda
\begin{pmatrix}
x_1+x_3\\
-x_2+2x_3\\
3x_1-x_2+x_3
\end{pmatrix}
=
\lambda \varphi(x).
\]

Следовательно, оператор из пункта а) является линейным.

\bigskip

\textbf{б)} Проверим, является ли линейным оператор
\[
\varphi(x)=
\begin{pmatrix}
x_1+x_3\\
-x_2+2x_3\\
3x_1-x_2+1
\end{pmatrix}.
\]

Пусть
\[
x=(x_1,x_2,x_3)^T,\qquad y=(y_1,y_2,y_3)^T.
\]

Тогда
\[
x+y=(x_1+y_1,\;x_2+y_2,\;x_3+y_3)^T.
\]

\[
\varphi(x+y)=
\begin{pmatrix}
(x_1+y_1)+(x_3+y_3)\\
-(x_2+y_2)+2(x_3+y_3)\\
3(x_1+y_1)-(x_2+y_2)+1
\end{pmatrix}
\]

\[
=
\begin{pmatrix}
x_1+y_1+x_3+y_3\\
-x_2-y_2+2x_3+2y_3\\
3x_1+3y_1-x_2-y_2+1
\end{pmatrix}.
\]

С другой стороны,
\[
\varphi(x)+\varphi(y)=
\begin{pmatrix}
x_1+x_3\\
-x_2+2x_3\\
3x_1-x_2+1
\end{pmatrix}
+
\begin{pmatrix}
y_1+y_3\\
-y_2+2y_3\\
3y_1-y_2+1
\end{pmatrix}
\]

\[
=
\begin{pmatrix}
(x_1+x_3)+(y_1+y_3)\\
(-x_2+2x_3)+(-y_2+2y_3)\\
(3x_1-x_2+1)+(3y_1-y_2+1)
\end{pmatrix}
\]

\[
=
\begin{pmatrix}
x_1+y_1+x_3+y_3\\
-x_2-y_2+2x_3+2y_3\\
3x_1+3y_1-x_2-y_2+2
\end{pmatrix}.
\]

Сравнивая результаты, видим, что
\[
\varphi(x+y)\neq \varphi(x)+\varphi(y),
\]
так как в третьей координате получаем
\[
3x_1+3y_1-x_2-y_2+1
\neq
3x_1+3y_1-x_2-y_2+2.
\]

Следовательно, аддитивность не выполняется, поэтому оператор не является линейным.

\bigskip

\textbf{в)} Проверим, является ли линейным оператор
\[
\varphi(x)=
\begin{pmatrix}
x_1+x_3\\
-x_2+2x_3\\
3x_1-x_2^2+x_3
\end{pmatrix}.
\]

Пусть
\[
x=(x_1,x_2,x_3)^T,\qquad y=(y_1,y_2,y_3)^T.
\]

Тогда
\[
x+y=(x_1+y_1,\;x_2+y_2,\;x_3+y_3)^T.
\]

\[
\varphi(x+y)=
\begin{pmatrix}
(x_1+y_1)+(x_3+y_3)\\
-(x_2+y_2)+2(x_3+y_3)\\
3(x_1+y_1)-(x_2+y_2)^2+(x_3+y_3)
\end{pmatrix}
\]

\[
=
\begin{pmatrix}
x_1+y_1+x_3+y_3\\
-x_2-y_2+2x_3+2y_3\\
3x_1+3y_1-(x_2+y_2)^2+x_3+y_3
\end{pmatrix}.
\]

С другой стороны,
\[
\varphi(x)+\varphi(y)=
\begin{pmatrix}
x_1+x_3\\
-x_2+2x_3\\
3x_1-x_2^2+x_3
\end{pmatrix}
+
\begin{pmatrix}
y_1+y_3\\
-y_2+2y_3\\
3y_1-y_2^2+y_3
\end{pmatrix}
\]

\[
=
\begin{pmatrix}
(x_1+x_3)+(y_1+y_3)\\
(-x_2+2x_3)+(-y_2+2y_3)\\
(3x_1-x_2^2+x_3)+(3y_1-y_2^2+y_3)
\end{pmatrix}
\]

\[
=
\begin{pmatrix}
x_1+y_1+x_3+y_3\\
-x_2-y_2+2x_3+2y_3\\
3x_1+3y_1-x_2^2-y_2^2+x_3+y_3
\end{pmatrix}.
\]

Так как
\[
(x_2+y_2)^2=x_2^2+2x_2y_2+y_2^2,
\]
то
\[
3x_1+3y_1-(x_2+y_2)^2+x_3+y_3
=
3x_1+3y_1-x_2^2-2x_2y_2-y_2^2+x_3+y_3.
\]

Следовательно, третья координата вектора $\varphi(x+y)$ равна
\[
3x_1+3y_1-x_2^2-2x_2y_2-y_2^2+x_3+y_3,
\]
а третья координата вектора $\varphi(x)+\varphi(y)$ равна
\[
3x_1+3y_1-x_2^2-y_2^2+x_3+y_3.
\]

Сравнивая результаты, видим, что
\[
\varphi(x+y)\neq \varphi(x)+\varphi(y),
\]
так как появляется дополнительное слагаемое $-2x_2y_2$.

Следовательно, аддитивность не выполняется, поэтому оператор не является линейным.

\bigskip

Итак, линейным является только оператор из пункта \textbf{а)}.

\bigskip

\textbf{2. Найдём матрицу оператора из пункта а) в стандартном базисе $e_1,e_2,e_3$ и в базисе $e_1',e_2',e_3'$, где}
\[
e_1' = e_1 - e_2 + e_3,
\qquad
e_2' = -e_1 + e_2 - 2e_3,
\qquad
e_3' = -e_1 + 2e_2 + e_3.
\]

\bigskip

\textbf{2.1. Матрица оператора в стандартном базисе.}

Найдём образы базисных векторов:
\[
\varphi(e_1)=\varphi
\begin{pmatrix}
1\\0\\0
\end{pmatrix}
=
\begin{pmatrix}
1+0\\
-0+0\\
3\cdot1-0+0
\end{pmatrix}
=
\begin{pmatrix}
1\\0\\3
\end{pmatrix},
\]
\[
\varphi(e_2)=\varphi
\begin{pmatrix}
0\\1\\0
\end{pmatrix}
=
\begin{pmatrix}
0+0\\
-1+0\\
0-1+0
\end{pmatrix}
=
\begin{pmatrix}
0\\-1\\-1
\end{pmatrix},
\]
\[
\varphi(e_3)=\varphi
\begin{pmatrix}
0\\0\\1
\end{pmatrix}
=
\begin{pmatrix}
0+1\\
-0+2\cdot1\\
0-0+1
\end{pmatrix}
=
\begin{pmatrix}
1\\2\\1
\end{pmatrix}.
\]

Следовательно, матрица оператора в стандартном базисе:
\[
A=
\begin{pmatrix}
1 & 0 & 1\\
0 & -1 & 2\\
3 & -1 & 1
\end{pmatrix}.
\]

\bigskip

\textbf{2.2. Матрица перехода к базису $e_1',e_2',e_3'$.}

Координаты новых базисных векторов в старом базисе:
\[
e_1'=
\begin{pmatrix}
1\\-1\\1
\end{pmatrix},
\qquad
e_2'=
\begin{pmatrix}
-1\\1\\-2
\end{pmatrix},
\qquad
e_3'=
\begin{pmatrix}
-1\\2\\1
\end{pmatrix}.
\]

Поэтому матрица перехода:
\[
T=
\begin{pmatrix}
1 & -1 & -1\\
-1 & 1 & 2\\
1 & -2 & 1
\end{pmatrix}.
\]

\bigskip

\textbf{2.3. Матрица оператора в новом базисе.}

Используем формулу
\[
A' = T^{-1}AT.
\]

Сначала найдём обратную матрицу $T^{-1}$.

\[
T=
\begin{pmatrix}
1 & -1 & -1\\
-1 & 1 & 2\\
1 & -2 & 1
\end{pmatrix}.
\]

Для этого воспользуемся методом присоединения единичной матрицы:
\[
\left(
\begin{array}{ccc|ccc}
1 & -1 & -1 & 1 & 0 & 0\\
-1 & 1 & 2 & 0 & 1 & 0\\
1 & -2 & 1 & 0 & 0 & 1
\end{array}
\right)
\sim
\left(
\begin{array}{ccc|ccc}
1 & -1 & -1 & 1 & 0 & 0\\
0 & 0 & 1 & 1 & 1 & 0\\
0 & -1 & 2 & -1 & 0 & 1
\end{array}
\right)
\]

\[
\sim
\left(
\begin{array}{ccc|ccc}
1 & -1 & -1 & 1 & 0 & 0\\
0 & -1 & 2 & -1 & 0 & 1\\
0 & 0 & 1 & 1 & 1 & 0
\end{array}
\right)
\sim
\left(
\begin{array}{ccc|ccc}
1 & -1 & -1 & 1 & 0 & 0\\
0 & 1 & -2 & 1 & 0 & -1\\
0 & 0 & 1 & 1 & 1 & 0
\end{array}
\right)
\]

\[
\sim
\left(
\begin{array}{ccc|ccc}
1 & -1 & 0 & 2 & 1 & 0\\
0 & 1 & 0 & 3 & 2 & -1\\
0 & 0 & 1 & 1 & 1 & 0
\end{array}
\right)
\sim
\left(
\begin{array}{ccc|ccc}
1 & 0 & 0 & 5 & 3 & -1\\
0 & 1 & 0 & 3 & 2 & -1\\
0 & 0 & 1 & 1 & 1 & 0
\end{array}
\right).
\]

Следовательно,
\[
T^{-1}=
\begin{pmatrix}
5 & 3 & -1\\
3 & 2 & -1\\
1 & 1 & 0
\end{pmatrix}.
\]

Теперь найдём $A'$:
\[
A'=T^{-1}AT.
\]

Сначала вычислим произведение $AT$:
\[
AT=
\begin{pmatrix}
1 & 0 & 1\\
0 & -1 & 2\\
3 & -1 & 1
\end{pmatrix}
\begin{pmatrix}
1 & -1 & -1\\
-1 & 1 & 2\\
1 & -2 & 1
\end{pmatrix}.
\]

\[
AT=
\begin{pmatrix}
\substack{1\cdot1+0\cdot(-1)\\+1\cdot1} &
\substack{1\cdot(-1)+0\cdot1\\+1\cdot(-2)} &
\substack{1\cdot(-1)+0\cdot2\\+1\cdot1} \\

\substack{0\cdot1+(-1)\cdot(-1)\\+2\cdot1} &
\substack{0\cdot(-1)+(-1)\cdot1\\+2\cdot(-2)} &
\substack{0\cdot(-1)+(-1)\cdot2\\+2\cdot1} \\

\substack{3\cdot1+(-1)\cdot(-1)\\+1\cdot1} &
\substack{3\cdot(-1)+(-1)\cdot1\\+1\cdot(-2)} &
\substack{3\cdot(-1)+(-1)\cdot2\\+1\cdot1}
\end{pmatrix}
\]

\[
=
\begin{pmatrix}
2 & -3 & 0\\
3 & -5 & 0\\
5 & -6 & -4
\end{pmatrix}.
\]

Теперь умножим $T^{-1}$ на $AT$:
\[
A'=T^{-1}(AT)=
\begin{pmatrix}
5 & 3 & -1\\
3 & 2 & -1\\
1 & 1 & 0
\end{pmatrix}
\begin{pmatrix}
2 & -3 & 0\\
3 & -5 & 0\\
5 & -6 & -4
\end{pmatrix}.
\]

\[
A'=
\begin{pmatrix}
\substack{5\cdot2+3\cdot3\\+(-1)\cdot5} &
\substack{5\cdot(-3)+3\cdot(-5)\\+(-1)\cdot(-6)} &
\substack{5\cdot0+3\cdot0\\+(-1)\cdot(-4)} \\

\substack{3\cdot2+2\cdot3\\+(-1)\cdot5} &
\substack{3\cdot(-3)+2\cdot(-5)\\+(-1)\cdot(-6)} &
\substack{3\cdot0+2\cdot0\\+(-1)\cdot(-4)} \\

\substack{1\cdot2+1\cdot3\\+0\cdot5} &
\substack{1\cdot(-3)+1\cdot(-5)\\+0\cdot(-6)} &
\substack{1\cdot0+1\cdot0\\+0\cdot(-4)}
\end{pmatrix}
\]

\[
=
\begin{pmatrix}
14 & -24 & 4\\
7 & -13 & 4\\
5 & -8 & 0
\end{pmatrix}.
\]

Следовательно, матрица оператора в базисе $e_1',e_2',e_3'$ равна
\[
A'=
\begin{pmatrix}
14 & -24 & 4\\
7 & -13 & 4\\
5 & -8 & 0
\end{pmatrix}.
\]
